\section{Setup}
\label{sec:setup}

This section describes how to build \vannak{}, bring up the supporting
infrastructure, and verify a working installation.

\subsection{Prerequisites}

\begin{description}
\item[Rust] stable toolchain with edition 2024 support. The Docker build uses
  \texttt{rust:1-bookworm}.
\item[Docker] with the Compose plugin, for PostgreSQL and the \raft{} cluster.
\item[protobuf-compiler] required by the \graft{} transport crates when
  building the \texttt{node} feature. On Debian/Ubuntu install with
  \shellcmd{apt-get install protobuf-compiler}.
\item[dnsutils] (the \shellcmd{dig} command) is needed at runtime by
  \texttt{vannak-node} for DNS SRV peer discovery. The runtime image installs
  it automatically.
\end{description}

\subsection{Sibling Repositories}

\vannak{} depends on two sibling repositories via local path dependencies. They
must be checked out next to the \vannak{} repository:

\begin{lstlisting}[numbers=none]
parent/
  vannak/   (this repository)
  ipto/     (Ipto metadata system; ../ipto/implementations/rust)
  raft/     (graft Raft implementation; ../raft/graft-rust)
\end{lstlisting}

The \texttt{ipto-writer} feature uses \filepath{../ipto/implementations/rust},
and the \texttt{raft} and \texttt{node} features use the crates under
\filepath{../raft/graft-rust}. The PostgreSQL Compose file mounts the \ipto{}
SQL schema from \filepath{../ipto/shared/db/postgresql}.

\subsection{Building}

The dependency-free core builds and tests with no feature flags:

\begin{lstlisting}[numbers=none]
cargo build
cargo test
cargo fmt --check
cargo clippy --all-targets -- -D warnings
\end{lstlisting}

Building with feature flags enables the optional adapters:

\begin{lstlisting}[numbers=none]
cargo build --features ipto-writer
cargo build --features raft
cargo build --features node --bin vannak-node
\end{lstlisting}

\begin{vnote}
The \texttt{node} feature implies \texttt{raft}. The \texttt{vannak-node}
binary is only compiled when the \texttt{node} feature is active.
\end{vnote}

\subsection{Docker Compose for PostgreSQL}

The file \filepath{docker-compose.test.yml} starts a PostgreSQL 16 instance
preloaded with the \ipto{} schema, stored procedures, and boot data. The
integration test connects to this instance.

\begin{lstlisting}[numbers=none]
docker compose -f docker-compose.test.yml up -d
VANNAK_PG_INTEGRATION=1 cargo test --features ipto-writer --test vannak_integration
docker compose -f docker-compose.test.yml down -v
\end{lstlisting}

The Compose service sets the environment variables \texttt{POSTGRES\_USER},
\texttt{POSTGRES\_PASSWORD}, and \texttt{POSTGRES\_DB} (all \texttt{repo}), and
exposes port 5432 (override with \texttt{VANNAK\_PG\_PORT}). Three SQL files are
mounted into \filepath{/docker-entrypoint-initdb.d} in order:

\begin{table}[htbp]
\centering
\begin{tabular}{p{0.24\textwidth}p{0.62\textwidth}}
\hline
Mounted file & Purpose \\
\hline
\filepath{00_schema.sql} & \ipto{} table schema \\
\filepath{01_procedures.sql} & stored procedures \\
\filepath{02_boot.sql} & boot data (tenant id 1 = \texttt{SCRATCH}) \\
\hline
\end{tabular}
\caption{PostgreSQL initialization files for the test backend.}
\end{table}

\subsection{Database Initialization}

PostgreSQL runs the files under \filepath{/docker-entrypoint-initdb.d} once, on
first start of an empty data volume. The \filepath{boot.sql} script seeds tenant
id 1 (\texttt{SCRATCH}), which \texttt{IptoRepoWriter} uses by default. To
re-run initialization, tear the volume down with \shellcmd{docker compose -f docker-compose.test.yml down -v}.

\subsection{Docker Compose for the Raft Cluster}

The file \filepath{docker-compose.cluster.yml} starts a three-node \vannak{}
\raft{} cluster with DNS-based auto-discovery. A CoreDNS container serves static
SRV records so the nodes find each other at startup.

\begin{lstlisting}[numbers=none]
docker compose -f docker-compose.cluster.yml build
docker compose -f docker-compose.cluster.yml up -d
docker compose -f docker-compose.cluster.yml down -v
\end{lstlisting}

The topology assigns static addresses on the \texttt{raft-net} bridge network
(subnet \texttt{10.77.0.0/24}):

\begin{table}[htbp]
\centering
\begin{tabular}{lll}
\hline
Container & Address & Role \\
\hline
\texttt{vannak-dns} & \texttt{10.77.0.10} & CoreDNS server \\
\texttt{vannak-1} & \texttt{10.77.0.11} & \raft{} node \\
\texttt{vannak-2} & \texttt{10.77.0.12} & \raft{} node \\
\texttt{vannak-3} & \texttt{10.77.0.13} & \raft{} node \\
\hline
\end{tabular}
\caption{Cluster Compose topology and static addresses.}
\end{table}

Each node listens on port 10081 and starts with
\shellcmd{--srv _vannak._tcp.vannak.local 0.0.0.0 10081 vannak-N /data}, where
the SRV service name resolves to the three node A-records through CoreDNS.

\subsection{Integration Test Scripts}

The script \filepath{scripts/run-integration-test.sh} orchestrates both
infrastructure stacks.

\begin{table}[htbp]
\centering
\begin{tabular}{p{0.24\textwidth}p{0.62\textwidth}}
\hline
Invocation & Behavior \\
\hline
\shellcmd{./scripts/run-integration-test.sh} & PostgreSQL + \ipto{} writer flow \\
\shellcmd{./scripts/run-integration-test.sh --cluster} & also start the \raft{} cluster \\
\shellcmd{./scripts/run-integration-test.sh --kafka} & also start Redpanda and run the Kafka-to-\sitas{} smoke test \\
\shellcmd{./scripts/run-integration-test.sh --kafka-only} & run only the Kafka-to-\sitas{} smoke test \\
\shellcmd{./scripts/run-integration-test.sh --no-build} & skip the build step \\
\shellcmd{./scripts/run-integration-test.sh --down} & tear down all stacks \\
\hline
\end{tabular}
\caption{Integration test runner modes.}
\end{table}

The script starts PostgreSQL, runs the \texttt{vannak\_integration} test with
\texttt{VANNAK\_PG\_INTEGRATION=1}, optionally starts Redpanda and runs
\texttt{kafka\_integration} with \texttt{VANNAK\_KAFKA\_INTEGRATION=1}, and,
when \texttt{--cluster} is given, builds the node image, brings up the three
nodes, waits for leader election, and probes each node.

\subsection{Verification}

\begin{enumerate}
\item \shellcmd{cargo test} succeeds for the dependency-free core.
\item With PostgreSQL up and \texttt{VANNAK\_PG\_INTEGRATION=1}, the
  \texttt{vannak\_integration} test passes the full ingest, index, outbox, and
  \ipto{} writer flow.
\item With Redpanda up and \texttt{VANNAK\_KAFKA\_INTEGRATION=1}, the
  \texttt{kafka\_integration} test produces a Durga JSON process event to Kafka
  and verifies that the \texttt{kafka-client} consumer feeds \sitas{} workers.
\item For the cluster, \shellcmd{docker compose -f docker-compose.cluster.yml ps}
  shows four running containers (DNS plus three nodes), and
  \shellcmd{vannak-node probe localhost 10081} returns a cluster summary.
\end{enumerate}

\subsection{Feature Flag Reference}

\begin{table}[htbp]
\centering
\begin{tabular}{ll}
\hline
Feature & Pulled-in dependencies \\
\hline
\texttt{ipto-writer} & \texttt{ipto\_rust}, \texttt{serde\_json}, \texttt{uuid} \\
\texttt{raft} & \texttt{graft-core}, \texttt{serde}, \texttt{serde\_json} \\
\texttt{sitas-runtime} & local \texttt{sitas} crate \\
\texttt{kafka-client} & \texttt{sitas-runtime} plus \texttt{rdkafka} \\
\texttt{node} & \texttt{raft} plus \texttt{graft-runtime}, \texttt{graft-storage}, \texttt{graft-transport}, \texttt{graft-proto}, \texttt{tokio}, \texttt{parking\_lot}, \texttt{bytes} \\
\hline
\end{tabular}
\caption{Feature flags and the dependencies they enable.}
\end{table}

\endinput
